% 本文件由 LaTeX 排版 Agent 根据有效 translation.json 自动生成。
% author=John Chrysostom; work_id=1905; title=关于圣依纳爵的讲道

\chapter{关于圣依纳爵的讲道}

% structure: p0001 source=h1 target=omitted reason=page-title-already-rendered-by-parent

\Emph{颂词。论神圣殉道者、背负天主者、伟大的\Place{安提约基雅}大主教圣\Person{依纳爵}，他被解送至\Place{罗马}，在那里殉道，然后又从那里被送回\Place{安提约基雅}。}\par

1. \Emph{丰盛的}而豪华的款待者经常不断地设宴款待，既为了显示自己的财富，也为了向相识表示善意。同样，圣神的恩宠提供其大能的证明，并向天主的友人显示极大的善意，接连不断地在我们面前摆上殉道者的筵席。例如最近，一位十分年轻且未婚的少女，真福殉道者\Person{佩拉吉亚}，以极大的喜乐款待了我们。今天，这位真福而高贵的殉道者\Person{依纳爵}又接续了她的庆节。人虽不同，筵席却是一样；角力方式各异，冠冕却是一样；竞赛多种多样，奖品却是同一。因为在异教人的竞赛中，由于任务是体力的，合理只有男子参加。但在这里，因为竞赛完全关乎灵魂，竞赛场向两性开放，每种性别都设有剧场。不仅男子单独脱衣，以免女子以本性软弱为托词，似乎有合理的借口；也不仅女子独自像男子般勇敢，以免男子蒙羞；而是此侧与彼侧，许多人被宣称为得胜者，并获得冠冕，好叫你借事迹本身知道，在基督耶稣内，不再分男女性别，\BibleRef{迦 3:28}也不分性别、身体软弱、年龄或任何这类事物，都不能阻碍那些在信仰赛程中奔跑的人；只要在我们灵魂中建立高尚的预备、热切的心志与炽热燃起的敬畏天主之情。为此，少女、妇女、男子、青年、老人、奴隶、自由人，以及每个身份、每个年龄、每个性别，都为此竞赛脱下衣服，且丝毫不受伤害，因为他们带来了高尚的目的来角力。那么，时节已召唤我们谈论这位圣人的伟大事迹。但我们的思绪纷乱而困惑，不知先说哪样，后说哪样，再说哪样，如此众多要求颂扬的事从四面围绕我们；我们所经历的，正如有人走进一片草地，看见许多玫瑰丛、许多紫罗兰，以及大量的百合花，还有其他多种多样的春花，便疑惑应先看哪一朵，后看哪一朵，因为他所见的每一朵都邀请他将目光投向自己。同样，我们来到这为\Person{依纳爵}圣迹所开辟的属灵草地，所见并非春天的花朵，而是此人灵魂中圣神那多样而丰富的果实，便混乱困惑，不知应先关注哪一样，因为我们所见到的每件事都吸引我们离开邻近的事，并诱使灵魂之眼观看其自身的美。请看，他尊贵地治理了我们中间的教会，其谨慎正是基督所愿。因为基督所宣告的主教职务的最高标准与准则，此人都以行为彰显了。因他听到基督说：\BibleQuote{善牧为羊舍掉自己的生命}\BibleRef{若 10:11}，他便以全部勇气为羊舍了命。\par

他真正与宗徒们交往，并饮于属灵的泉源。那么，这样一个人，曾受培育，处处与他们交往，分享可说的与不可说的真理，又被他们视为堪当如此崇高职位，他会是怎样的人呢？时机再次来临，要求勇气；一个蔑视一切现世事物、为天主之爱炽热、以未见之事重于可见之事的灵魂，他放下肉身如脱去一件衣服般轻松。那么，我们应先说什么呢？是他自始至终所证实的宗徒教导，还是他对现世生命的漠视，或是他治理教会时所施展的德行之严谨？我们应先想起哪一样？是殉道者、主教、还是宗徒？因为圣神的恩宠编织了一顶三重冠冕，如此绑在他圣洁的头上，不，是多重的冠冕。因为若有人仔细考量，他会发现每一顶冠冕都为我们绽放出其他冠冕。\par

2. 若你愿意，我们先来赞美他的主教之职。这似乎只是一顶冠冕吗？那么我们且用言语展开它，你将会看到从中产生了两个、三个甚至更多。因为我不仅惊叹此人本身堪当如此伟大的职务，更惊叹他从那些圣者手中获得了这一职务，而且有福的宗徒们的手触摸了他神圣的头。因为说他值得赞美，这并非小事，也不是因为他从高处赢得了更大的恩宠，也不仅是因为他们使圣神更充沛的德能临于他身上，而是因为他们为他作证，证明人所拥有的一切美德都在他内。现在我来告诉你们这是如何的。\Person{保禄}曾写信给\Person{弟铎}——当我说\Person{保禄}时，我不仅是指他一人，也包括\Person{伯多禄}、\Person{雅各伯}、\Person{若望}以及他们全体；因为正如在一架竖琴中，琴弦各不相同，但和谐却是一个，同样，在宗徒的团体中，人物各不相同，但教导却是一个，因为工匠是同一个，我是指那推动他们灵魂的圣神，\Person{保禄}显示这一点时说道：\BibleQuote{所以，不拘是我，或是他们，我们都这样宣讲}\BibleRef{格前15:11}。于是，此人写信给\Person{弟铎}，指示主教应当是怎样的人，他说：\BibleQuote{因为作主教的，必须无可指责，如同天主的管家；不自高自大，不轻易发怒，不好酒，不打人，不贪图不义之财；却要好客，喜爱良善，端庄，正义，圣洁，节制，持守那合乎教理、可信赖的道理，使他既能以健全的道理劝勉人，又能驳倒反对的人。}\BibleRef{铎1:7}。并且又写信给\Person{弟茂德}，论及同一主题时，他大致这样说：若有人渴望主教的职务，他向往的是善工。所以，主教必须无可指责，只有一个妻子的丈夫，节制，庄重，有序，好客，善于教导，不好酒，不打人，但温和，不好争辩，不贪爱钱财。你们看见他要求主教有何等严格的美德吗？因为正如一位最出色的画家，调配了许多色彩，若要描绘一幅君王形体的原稿，就会一丝不苟地工作，使所有临摹和照着它绘画的人都能得到准确的肖像；同样，有福的\Person{保禄}，正如在描绘一幅君王的肖像，提供其原稿，调配了各种美德的色彩，完整地画出了主教职务的特征，以便每一个登上这尊位的人，观看此像，就能以同样严格的态度管理自己的事务。\par

大胆地说，\Person{依纳爵}在自己灵魂中精确地印刻了这一切：他是无可指责、无可指摘的，不固执己见，不轻易发怒，不酗酒，不打人，却温和，不好争辩，不贪爱钱财，正义，圣洁，节制，持守那合乎教理、可信赖的道理，清醒，庄重，有序，以及\Person{保禄}所要求的一切其他品德。有人说：\Quote{这有什么证据呢？}那些说了这些话的人祝圣了他，而那些如此严格地建议别人考核那些将要登上这职务宝座的人，他们自己不会轻率地行这事。但若他们没有看到这所有美德植根于这位殉道者的灵魂中，就不会将这职务托付给他。因为他们清楚地知道，那些粗心大意、随意进行此类祝圣的人面临着多大的危险。\Person{保禄}再次向同一\Person{弟茂德}指明这事时写道：\BibleQuote{不可轻易给人覆手，也不可在别人的罪上有份。}\BibleRef{弟前5:22}。你怎么说？别人犯了罪，我也要分担他的责备和惩罚吗？是的，他说，那认可邪恶的人；正如有人将一把锋利的剑交到一个狂暴的疯子手中，那疯子用剑杀人，交剑的人要承担责任；同样，若有人将这职务所赋予的权力交给一个生活在邪恶中的人，就是把那人一切罪恶和狂妄的火焰引到自己头上。因为那提供根源的人，就是产生周围一切后果的原因。你们看见在此期间主教的职务如何出现了一顶双重的冠冕，以及那些祝圣他的人所具有的尊贵如何使这职务更加光荣，为他身上的每一美德作了见证？\par

3. 你愿意我也向你揭示从同一件事而来的另一顶荣冠吗？让我们思考他获得这一尊荣的时刻。因为如今管理教会与那时管理教会并不相同，如同行走在一条经过许多行人而踩实、预备好的路上，与行走在一条刚刚开辟、满是车辙、石块、野兽、且从未有旅行者走过的路上并不相同。因为如今，因天主的恩宠，主教们没有危险，而是四面深享平安，我们都得享宁静，因为虔诚的圣言已传到地极，我们的统治者严谨持守信德。但那时并无这些，无论人往何处看，都是悬崖和陷阱、战争、搏斗和危险；统治者、君王、人民、城市、民族，以及国内国外的人，都为信徒设下陷阱。这并非唯一严重的事，还有许多信徒自身，因初次尝到陌生的教义，急需极大的宽容，且仍处于较为软弱的状态，时常动摇。这使教师们忧伤，其程度不亚于外来的争斗，甚至更甚。因为外来的争斗和阴谋，因着对未来赏报的希望，给他们带来许多喜乐。为此，宗徒们从公议会面前回来，因为被打而欢喜\BibleRef{宗 5:41}；保禄呼喊说：\Quote{我因所受的苦而喜乐}\BibleRef{哥 1:24}，他处处以磨难为荣。但家中的创伤和弟兄的失足不容他们喘息，反而如同最沉重的轭，持续压迫折磨他们灵魂的颈项。请听保禄虽在苦难中喜乐，却为这些事痛苦煎熬。他说：\Quote{有谁软弱，我不软弱呢？有谁跌倒，我不心焦呢？}\BibleRef{格后 11:29}又说：\Quote{我怕当我来的时候，见到你们不合我所希望，而你们见我也不合你们所希望。}\BibleRef{格后 12:20}稍后又说：\Quote{恐怕我再来的时候，我的天主使我在你们面前谦卑，我要为许多从前犯罪而不悔改、仍然行不洁、淫乱和放荡的人哀哭。}\BibleRef{格后 12:21}通篇可见，他为家中的成员流泪哀恸，为信徒时时恐惧战栗。正如我们敬佩舵手，不是在他能将船上的人安全渡岸、海面平静、船顺风而行之时，而是在深渊翻腾、波涛相争、乘客内讧、船中内外遭遇巨大风暴、而他能稳妥驾船之时；同样，我们应当惊讶并敬佩那些当时受托掌管教会的人，远胜于现今管理教会的人。当时内外有大战，信德的幼苗更为嫩弱，需要极多的照料；教会中的群众如同新生婴孩，亟需极大的远见和任何预备哺养它的灵魂所具备的最大智慧。为使你更清楚知道，那些受托管理教会的人当时配得何等大的荣冠，以及在门槛和开端承担此事、首先进入其中是多么大的工作和危险，我为你提出基督的见证，祂对这些事作出判断，并确认我所表达的见解。因祂看见许多人来到祂那里，并愿向宗徒们指出先知们比他们更劳苦，祂说：\Quote{别人劳苦，你们却进入他们的劳苦。}\BibleRef{若 4:38}然而宗徒们比先知们劳苦得多。但因他们首先播下虔诚的圣言，将人未受教导的灵魂赢向真理，大部分工作便归功于他们。因为在许多教师之后来教导，与亲自首先播下种子，绝不相同。因为已实践过、成为许多人习惯的事，容易接受；但如今初次听到的事，会扰乱听者的心思，给教师带来大量工作。至少这正是扰乱雅典听众的事，为此他们转离保禄，责备他说：\Quote{你把一些奇怪的事传到我们耳中。}\BibleRef{宗 17:20}因为若现今教会的监管给治理者带来许多辛劳和工作，那么想一想当时的工作是何等加倍、三倍、多倍，当时有危险、争斗、陷阱和持续的恐惧。无法用言语充分表达那些圣人当时所遭遇的困难，只有亲身经历的人才能知道。\par

4. 现在我要谈第四顶冠冕，是由这主教职务为我们产生的。那么这顶是什么呢？就是他受托管理我们本城的事实。因为管理一百人，甚至仅五十人，确实是劳累的事。但若承担如此巨大的城市，人口达二十万，你认为这是何等品德与智慧的证明呢？因为如同在军队的照管中，较明智的将军掌管更精锐、更庞大的军团，同样，在城市的管理中，较能干的统治者受托管理更大、人口更多的城市。无论如何，这座城市在天主眼中颇为重要，祂确实以祂所行的事迹显明了这一点。无论如何，整个世界的师傅\Person{伯多禄}，祂将天国的钥匙交在他手中，命令他行事并承受一切，祂命他在此长住。因此，在祂眼中，我们的城相当于整个世界。但是，既然我提到了\Person{伯多禄}，我看出有一顶由他编织的第五冠冕，那就是此人继承了他的职务。因为正如有人从根基中取出一块大石，必设法引入一块同等大小的石头，以免动摇整个建筑，使其更不稳固；同样，当\Person{伯多禄}即将离开此地时，圣神的恩宠引入了另一位与\Person{伯多禄}同等的导师，以免已完成的建筑因继任者的卑微而变得更不稳固。这样，我们算出了五顶冠冕：来自职务的重要性，来自祝圣者的尊贵，来自时代的艰难，来自城市的大小，来自将主教职务传给他的人的德行。编织了所有这些之后，本可以再谈第六、第七甚至更多；但为了不把全部时间花在讨论主教职务上而错过关于殉道者的细节，从这点起，让我们转向那场战斗。曾有一次，一场残酷的战争重新燃起，反对教会；好像一种最残酷的暴政覆盖大地，所有人都从广场中央被带走。他们并非被控任何可怕之事，而是因为脱离谬误，奔向虔诚；因为戒绝事奉魔鬼，因为认识真天主，并崇拜祂的独生子；而为了这些他们本应受冠冕、赞美和荣耀的事，他们却受惩罚，遭遇无数折磨——所有信奉信仰的人，尤其是那些管理教会的人。因为魔鬼狡猾，善于策划这种阴谋，预料到如果它除去牧者，就很容易驱散羊群。但那位在狡猾中捉住狡猾者的天主，愿意向它显示：不是人管理祂的教会，而是祂亲自到处牧养那些信祂的人。祂允许此事发生，为让魔鬼看到，当牧者被除去时，虔诚的事业并未失败，宣讲的道并未熄灭，反而增加了；藉着这些事实，它自己以及所有事奉它的人可以学到，我们的事并非出于人，而是我们教导的根源来自高处，来自天上；是天主亲自引领教会，与天主作战者绝不能赢得胜利。但是魔鬼不仅行了这恶事，还行了一件不小于此的恶。因为不仅在主教们所管辖的城市中，它任凭主教被屠杀；还把他们带到外邦领土上杀死；它这样做，既急于趁他们失去朋友时下手，又希望用旅途的劳苦削弱他们。它果然用此对待这位圣人。因为它把他从我们的城召到\Place{罗马}，使路程加倍，期望以路途遥远和日数众多来压抑他的心，却不知有\Person{耶稣}与他同行，作为旅伴和同被放逐者，在这漫长旅程中，他反而变得更坚强，给他同在的能力提供了更多证明，并更紧密地联合了各教会。因为沿途的城市从四面八方聚集，鼓励这位运动员，用许多供应使他快行，以祈祷和代祷分享他的战斗。当他们看到殉道者如此欣然奔赴死亡，如一个被召往天国领域的人所应有的那样，并且透过事实本身，透过那位高尚者的欣然与喜乐，他们获得了莫大的安慰，知道他奔赴的并非死亡，而是一种长途旅行，是从此世的迁移，是升天；他离开时，在每个城市用言语和行为教导这些事，正如发生在犹太人身上的情形：当他们捆绑\Person{保禄}，送他到\Place{罗马}时，他们以为是在送他赴死，实际上却是在给住在那里的犹太人送一位导师。这事果然在\Person{依纳爵}身上更大程度地发生了。因为他不单对居住在\Place{罗马}的人，而且对所有位于中途的城市，都成了一位奇妙的导师，劝他们轻视现世生活，视所见之物为无，爱慕将来的事物，仰望天堂，完全不把现世任何恐怖放在心上。因为他以行为教导他们这些事，甚至更多；他一路前行教导他们，如同太阳从东方升起，急速趋向西方。但他比太阳更辉煌，因为太阳习惯于在高处运行，带来物质的光；但\Person{依纳爵}在下方照耀，将教理的精神之光赐予人的灵魂。那光一进入西方地区，就被隐藏，立刻带来黑夜；但这光进入西方时，在那里照耀得更加辉煌，给沿路所有地方带来最大的益处。当他到达那城时，甚至\Emph{那城}也被他教导了基督信仰的智慧。因此，天主允许他在那里结束生命，以使此人的死亡对\Place{罗马}所有居民都有教诲。因为\Emph{我们}藉天主的恩宠，已扎根于信德，不再需要证据。但住在\Place{罗马}的人，因为那里有不敬神的大罪，需要更多帮助。因此，\Person{伯多禄}、\Person{保禄}和后来的\Person{依纳爵}都在那里被杀：部分是为了用自己的血净化那座被偶像之血玷污的城市，部分是为了用他们的行为提供被钉\Person{基督}复活的证据，说服\Place{罗马}居民：他们不会如此愉快地蔑视现世生命，除非他们坚信自己即将升到被钉的\Person{耶稣}那里，在天上看见祂。因为实际上，被杀的\Person{基督}在死后显示出如此大的能力，以至于说服活着的人为了宣认祂而蔑视国家、家园、朋友、相识和生命本身，并且选择鞭打、危险和死亡来代替现世的快乐，这是复活的最大证明。因为这不是任何死人的成就，也不是留在坟墓中的人，而是复活并活着的人。因为你怎么解释：当祂活着时，所有与祂同行的宗徒因害怕而软弱，出卖他们的老师，逃跑离开；但当祂死后，不仅\Person{伯多禄}和\Person{保禄}，甚至\Person{依纳爵}——他从未见过祂，也未享受过与祂的陪伴——竟表现出如此的热忱，甘愿为祂舍命？\par

5. 为此，为使所有居住在\Place{罗马}的人得知这些事是真实的，天主准许那圣人就在那里得以成全；我愿从他死亡的方式本身来证实这一点。他并非在城墙外的地牢中，也非在法庭或某个角落接受定罪判决，而是在剧场中央，在全城百姓都坐在他上方时，他承受了这种形式的殉道——野兽被放出来攻击他——为在众人眼前竖立他对魔鬼的胜利，并使所有观众效法他的战斗。他不仅死得如此高贵，而且死得甚至带着喜乐。因为他看待野兽时并非如同即将与生命分离，而是如同被召往更美好、更属灵的生命一般，满心欢喜。这从何显明呢？从他临终时所说的话：当他听说等待他的是这种刑罚时，他说：\Quote{愿获得喜乐，}他说，\Quote{从这些野兽中。}因为爱德如此深厚：他们甘愿为所爱者承受任何痛苦，并且当所遭遇的苦难超乎寻常时，反倒感到心愿满足。这正发生在\Person{依纳爵}身上。因为他不仅藉着死亡，也藉着心甘情愿，努力效法宗徒们；他听说宗徒们受鞭打后欢喜离去，便也愿效法他的老师，不只在死亡上，也在喜乐上。为此他说：\Quote{愿在你们的野兽中获得喜乐。}他认为野兽之口远比暴君的舌头温和得多；这十分合理。因为暴君的舌头邀请他前往\OriginalTerm{革赫纳}{Gehenna}，而野兽之口却护送他进入天国。因此，当他在那里结束了生命——更确切地说，当他升入天堂时——他从此戴着冠冕离去。这也是出于天主的安排：天主把他归还给我们，并将这位殉道者分赐给各城。因为那座城领受了他滴落的鲜血，而你们却因他的遗骸而得尊荣；你们曾享有他的主教职，他们则享有他的殉道。他们看见他在战斗中得胜并加冕，而你们却永远拥有他。天主将他从你们身边带走片刻，却以更大的荣耀将他再次赐予你们。正如那些借钱的人连本带利归还所收受的，同样，天主使用了你们这份珍贵的宝藏片刻，并在那座城展示后，以更灿烂的光辉将它归还给你们。你们送去一位主教，却领回一位殉道者；你们以祈祷送他出发，却以冠冕迎接他归来；不仅你们，还有沿途各城。你们想，当他们看见他的遗骸被送回时，他们是如何行事的？何等喜悦！他们何等欢欣！他们从四面围着这位戴冠者，欢呼喝彩！正如一位高贵的运动员，摔倒所有对手后，带着灿烂的荣耀走出竞技场，观众迎接他，不让他脚沾地，将他扛在肩上送回家，并报以无数赞美：同样，各城依次从\Place{罗马}接下这位圣人，将他扛在肩上一直送到本城，以赞美护送这位戴冠者，歌唱颂扬这位冠军；嘲笑魔鬼，因为它的诡计反成了自己的损失，它想加害殉道者的事反而对他有利。他当时确实造福并鼓励了所有城市；并且从那时至今，他一直丰裕着这座城市；如同某种永久的宝藏，每日取用却永不枯竭，使所有分沾的人都更加昌盛；同样，这位蒙福的\Person{依纳爵}也使来到他面前的人充满祝福、勇气、高尚的精神和极大的胆量，然后送他们回家。\par

因此，不但今天，而且每天让我们到他那里去，从他身上采摘属灵的果实。因为带着信德来到此地的人，确实能够收获许多美善的果实。不仅圣人的身体，就连他们的坟墓也充满了属灵的恩宠。因为在厄里叟的事上曾发生这样的事：一具尸体碰到他的坟墓，就挣断了死亡的束缚，重新回到生命。\BibleRef{列下13:21} 如今恩宠更为丰沛，圣神的能力更为强大，带着信德触摸坟墓的人岂不更能获得大能？为此，天主允许我们保有圣人们的遗骸，愿意藉此把我们引向同样的热忱，并赐给我们一种港湾，以及对于不断袭来的种种邪恶的稳妥慰藉。因此我恳求你们众人：若有人沮丧、患病、受辱、或处于今生的任何其他困境，或深陷罪恶之中，让他带着信德来到这里，他必将放下这一切，带着极大的喜乐回去，因为仅凭瞻仰，良心得享轻快。不仅如此，不仅受苦者应来，就是那些处于喜乐、荣耀、权能、对天主有极大信心的人，也不要轻看这益处。因为他们来到这里，瞻仰这位圣人，藉着回忆此人的伟大事迹，能劝勉自己的灵魂保持节制，不致因这些伟业而自高自大。对于身处顺境之人，不为好运所飘然，而能以节制承受幸福，这绝非小事。因此，这宝藏对众人都有用：这安息之所适于落魄者，使他们脱离诱惑；适于幸运者，使他们的成功安稳；适于软弱者，使他们恢复健康；适于健康者，使他们不致陷入软弱。考虑到这一切，让我们把这样的时光看得高于一切愉悦、一切享乐，以便我们同时欢欣和获益，能成为这些圣人的同居者和同家人，藉着圣人们自己的祈祷，藉着我们的主耶稣基督的恩宠和慈爱，愿光荣借着祂归于圣父及圣神，现在及永远，世世无尽。阿们。\par

\SourceRecord{Translated by T.P. Brandram.；Nicene and Post-Nicene Fathers, First Series；Vol. 9.；Edited by Philip Schaff.；Buffalo, NY: Christian Literature Publishing Co.,；1889.；<http://www.newadvent.org/fathers/1905.htm>.}
